\documentclass[onecolumn, draftclsnofoot,10pt, compsoc]{IEEEtran}
\usepackage{graphicx}
\usepackage{url}
\usepackage{setspace}
\usepackage{listings}
\usepackage{xcolor}

\usepackage{geometry}
\geometry{textheight=9.5in, textwidth=7in}

% Listings settings for swift taken from chriseidhof GitHub
\lstdefinelanguage{swift}
{
  morekeywords={
    func,if,then,else,for,in,while,do,switch,case,default,where,break,continue,fallthrough,return,
    typealias,struct,class,enum,protocol,var,func,let,get,set,willSet,didSet,inout,init,deinit,extension,
    subscript,prefix,operator,infix,postfix,precedence,associativity,left,right,none,convenience,dynamic,
    final,lazy,mutating,nonmutating,optional,override,required,static,unowned,safe,weak,internal,
    private,public,is,as,self,unsafe,dynamicType,true,false,nil,Type,Protocol,
  },
  morecomment=[l]{//}, % l is for line comment
  morecomment=[s]{/*}{*/}, % s is for start and end delimiter
  morestring=[b]" % defines that strings are enclosed in double quotes
}

\definecolor{keyword}{HTML}{BA2CA3}
\definecolor{string}{HTML}{D12F1B}
\definecolor{comment}{HTML}{008400}

\lstset{
  language=swift,
  basicstyle=\ttfamily,
  showstringspaces=false, % lets spaces in strings appear as real spaces
  columns=fixed,
  keepspaces=true,
  keywordstyle=\color{keyword},
  stringstyle=\color{string},
  commentstyle=\color{comment},
}

% 1. Fill in these details
\def \CapstoneTeamName{Manovacumeter Team}
\def \CapstoneTeamNumber{ 20}
\def \GroupMemberOne{Cade Raichart}
\def \GroupMemberTwo{Daniel Kato}
\def \GroupMemberThree{Nathan Shepherd}
\def \CapstoneProjectName{ISS Barometer App }
\def \CapstoneSponsorCompany{NASA}
\def \CapstoneSponsorPerson{Don Pettit}

% 2. Uncomment the appropriate line below so that the document type works
\def \DocType{Progress Report}

\newcommand{\NameSigPair}[1]{\par
\makebox[2.75in][r]{#1} \hfil 	\makebox[3.25in]{\makebox[2.25in]{\hrulefill} \hfill		\makebox[.75in]{\hrulefill}}
\par\vspace{-12pt} \textit{\tiny\noindent
\makebox[2.75in]{} \hfil		\makebox[3.25in]{\makebox[2.25in][r]{Signature} \hfill	\makebox[.75in][r]{Date}}}}
% 3. If the document is not to be signed, uncomment the RENEWcommand below
%\renewcommand{\NameSigPair}[1]{#1}

%%%%%%%%%%%%%%%%%%%%%%%%%%%%%%%%%%%%%%%
\begin{document}
\begin{titlepage}
    \pagenumbering{gobble}
    \begin{singlespace}
    	\includegraphics[height=4cm]{coe_v_spot1}
        \hfill
        % 4. If you have a logo, use this includegraphics command to put it on the coversheet.
        %\includegraphics[height=4cm]{NASA}
        \par\vspace{.2in}
        \centering
        \scshape{
            \huge CS Capstone \DocType \par
            {\large\today}\par
            \vspace{.5in}
            \textbf{\Huge\CapstoneProjectName}\par
            \vfill
            {\large Prepared for}\par
            \Huge \CapstoneSponsorCompany\par
            \vspace{5pt}
            {\Large\NameSigPair{\CapstoneSponsorPerson}\par}
            {\large Prepared by }\par
            Group\CapstoneTeamNumber\par
            % 5. comment out the line below this one if you do not wish to name your team
            %\CapstoneTeamName\par
            \vspace{5pt}
            {\Large
                \NameSigPair{\GroupMemberOne}\par
                \NameSigPair{\GroupMemberTwo}\par
                \NameSigPair{\GroupMemberThree}\par
            }
            \vspace{20pt}
        }
        \begin{abstract}
        % 6. Fill in your abstract
        The International Space Station, the conglomerate work of multiple national space associations, including NASA, requires precision devices to monitor the many aspects of space life.
        One very important aspect for the success of life in near vacuum is sustained air pressure inside the modules.
        This air pressure is monitored by mechanical barometers, called Manovacometers, which must be visually monitored, and are no longer available.
        The \CapstoneProjectName seeks to provide a suitable replacement for the current barometer device aboard the International Space Station.
        This document will update the reader on our progress, from the last Fall term Progress report.
        Each member has written their progress and plans, as well as a main refresher on our project as a whole.
        The member sections include Responsibilities, Current Status, Remaining Tasks, Problems Faced, and additional relevant or interesting details.
        This document should provide a clear picture of where the project stands, halfway through the Winter term, and halfway through the allotted year.
        \end{abstract}
    \end{singlespace}
\end{titlepage}
\newpage
\pagenumbering{arabic}
\tableofcontents
% 7. uncomment this (if applicable). Consider adding a page break.
%\listoffigures
%\listoftables
\clearpage

% 8. now you write!
\section{Project Purpose}
The ISS Barometer Application will be used on the ISS to measure and track air pressure.
The crew members aboard the ISS are currently reliant on mechanical barometers to take critical readings during leaks, or rapid depressurization events.
These leaks, caused by orbital debris, must be found and patched or the crew would need to evacuate.
Our application will improve on the mechanical barometers by providing clear and accurate pressure readings through the iPads on-board barometer sensor.
Readouts and measurements will be presented in a easy-to-read fashion that will allow users to quickly gauge the important details.
In further iterations, we plan to implement contextualized audio data readouts by using the t-Res document and a simple Swift api \cite{probStat}.

\section{Project Goals}
As mentioned above the application we are making will provide ISS crew members with an accurate and efficient pressure reading.
Our required features follow that purpose and are listed below:
\begin{itemize}
\item[V1:] Data Screen with initial pressure, current pressure, start/stop button, time stamp, pressure change (measured in sec / mmHg)
\item[V1:] Graph Screen with plot pressure as a function of time
\item[V1:] Additional settings menu that includes plot scale, scroll functionality, orientation, decimal point, as well as other other units of measure for pressure
\item[V1:] All items will be optimized for efficient reading and clear interface
\item[V2:] Graph will have pinch to zoom functionality, scaling options (to be set in settings)
\item[V2:] Graph will be able to be exported as a csv, excel readable file
\item[V2:] Graph can be compared to t-Res table, to measure time remaining
\item[V2:] Audio data readouts that will provide specific information at regular intervals
\end{itemize}

Along with these specific features, the application that we create must be efficient and clear, as decided by the users.
The application will be used in challenging situations and must be clear and easy to understand, so that the data can be conveyed quickly and without confusion.
From the listed project goals, we have sufficiently finished with V1, and have started implementing portions of V2, while trying to finalize our V1 progress.


% These sections can be used by whoever wrote them, I'm not sure who did which ones though

% \section{Current State of the Project}
% % If anyone can think of a better way to state this, feel free to change the heading
% Currently we have an outline of what parts must be included in the app as well as the outline for when the first version of our app must be completed.
% We have also created our first prototype.
% This reads from the barometer and displays the current pressure in mm/Hg. It also calculates the $\frac{\Delta p}{\Delta t}$ and displays that alongside the current pressure.
% We have implemented the plot of pressure against time, and have had brief discussions with our client on design changes.

% \section{Project Problems}
% % Needs Rewritten
% We chose to use swift and Xcode because they allows for quick development and prototyping.
% This said we ran into many problems with that. For one Swift only works on Linux and MacOS.
% One of our group members did not have either at the time. They fixed this with installing a linux OS.
% Another problem was installing the swift libraries and connecting them with the Xcode application.
% Linux needed many packages to be downloaded for the library to work correctly, this took many hours of troubleshooting as it was not straightforward.
% Connecting swift to Xcode on the MacOS was another problem as well.
% Our last problem with swift and Xcode is we have no experience using either.
% We had and still have to continually work on improving our knowledge of these systems using online videos, guides and other resources.
%
% For the first time we all experienced what it was like working with someone who is in the workplace.
% We learned that our clients will not be getting back to us immediately and sometimes we have to wait for them reply.
% This means we have to work and finish our work and send it to him way earlier then we had planned originally.
% We also ran into the problems of finding the correct IEEE documentations.
% Sometimes we found the document quick and easy while other times it hours and many students combined to find them.


% We need to each write ~1000 words on our contributions
% Each of our sections could have interesting code, problems faced, issues managed
% I don't know really, but from the requirement this is the bulk of the writing

% Nathan- I've added sections for our bits. I think we can fill in Responsibilities with just the
%           issues we've taken for gitHub, I just copy and pasted them in including ones that I had a part in.
%
\section{Specific Contributions: Daniel Kato}
\subsection{Responsibilities}
\subsection{Current Status}
\subsection{Remaining Tasks}
\subsection{Problems Faced}
\subsection{Relevant and Interesting Details}

\section{Specific Contributions: Cade Raichart}
\subsection{Responsibilities}
\subsection{Current Status}
\subsection{Remaining Tasks}
\subsection{Problems Faced}
\subsection{Relevant and Interesting Details}

\section{Specific Contributions: Nathan Shepherd}

\subsection{Responsibilities}

From the technology document, I was tasked with the general projects of language choice, data storage, and color palette.
In those sections, I chose Swift, NSUserDefaults, and a self-defined blue theme.
Although these choices made significant impact on our design process, they didn't provide many facets of work.
The language choice affected all teammates, and was in fact already chosen at the start of the project.
The data storage is more complicated, and does have work involved with it; but that work was done by another teammate, as it aligned better with their past contributions.
The color palette turned out to be more restrained than expected, as the users have strict needs for readability and clarity.
This meant that our lettering and icons were required to be black on white background, for visibility, and only our banners and the application icons were altered.
Due to the technologies being vague and rather non-actionable, the team stuck to GitHub issues as our means of delegation.
As follows, these are the issues I took both full and partial responsibility for:

\begin{itemize}
  \item Adjust Y-Axis Default Zoom
  \item Random Debug Data
  \item Unit Picker Setting
  \item Rolling graph Setting
  \item Adding X-axis and Y-axis titles
\end{itemize}

\subsection{Current Status}

Currently our application is a fully functional prototype with near full coverage of our version one goals.
It properly records air pressure, and displays a graph, along with the requested current pressure, $\frac{\Delta p}{\Delta t}$, and $\frac{\Delta t}{\Delta p}$.
In addition we have implemented the following settings, and initialized them with assumed defaults:

\begin{lstlisting}
class Settings: Codable {
    var units = "mmHg"
    var sigFigs = 4
    var orientation = "Right"
    var slidingScale = false
    var windowSize = 50
}
\end{lstlisting}

These settings allow the user to change and save their preferences, making the application more interactive.
The units that can be chosen are mmHg, psi, kPa, and atm.
The changes then propagate to both the data view and the plot, making both scale appropriately.
The significant figure setting can be incremented anywhere between 0 and 10.
The orientation is a user setting because, as described in our requirements document, the iPad accelerometer doesn't work very well in the microgravity of the ISS.
To solve this the orientation is fixed, and can be changed by the user manually.
The sliding scale is a switch that will allow the user to choose between two graph functions: The first and default, is a scaling graph, that shrinks with more data; the second is a graph with a set window size, and will follow new data as it's returned.
The set window size option, if selected, will reveal a window size slider that allow the user to choose the size of the rolling window.
In addition to specific user settings, I have implemented a stream of debug data, when an on board barometer isn't available.
This will be removed before our final release, but it allows us to test the application just by running it on the Xcode emulator.

\subsection{Remaining Tasks}
The majority of remaining issues that we have to face are either UI adjustments, or version two features.
The UI adjustments include making the labels both larger and more readable, and adjusting the axis labels further (more info in the next section).
The client has also added some functionality to the version two release which include the following:

\begin{itemize}
  \item Dictation
  \item Screen lock
  \item T-Res Alarm
\end{itemize}

These issues will not be out of scope, and shouldn't be difficult to implement.
The dictation will include regular audio readouts of the pressure and pressure rates.
The interval and detail of these readouts have yet to be determined, but by implementing an internal text-to-speech Swift interface, the dictation should be quick to introduce and provide a hands free solution in the case of an emergency.

\subsection{Problems Faced}
\subsection{Relevant and Interesting Details}

% \section{Interesting Code}
% \begin{lstlisting}
%   func startDisplayingPressureData() {
%       self.altimeter.startRelativeAltitudeUpdates(to: OperationQueue.main, withHandler: {
%         data, error in
%           self.handlePressureReading(data: data!)
%       })
%   }
%
% \end{lstlisting}
%
% One of the main features of our application is its ability to read data from the iPad's built in barometer and display the data collected in different ways.
% The way that data is collected from the barometer is quite interesting.
% The process of reading data from the barometer requires a bit of processing power and time which can block other time sensitive tasks form being performed.
% As a result, the measurement is done asynchronously through an object called the \texttt{OperationQueue}.
% A method provided by the CoreMotion library called \texttt{startRelativeUpdates()} takes the \texttt{OperationQueue}'s \texttt{main} queue, and a callback function as arguement.
% The function then appends a task to the queue to perform in the background, and when the task is complete, the result is passed into the callback function for further processing.\\

\section{Retrospective}
\begin{center}
	\begin{tabular}{| l | p{0.3\linewidth} | p{0.3\linewidth} | p{0.3\linewidth} |}
		\hline
		& Positives & Deltas & Actions \\ \hline
    Winter Break & --  & -- & -- \\ \hline
    Week 1 & We made our GitHub setup with issues and a basic prototype
           & Started learning Swift and becoming familiar with the interface
           & Detailed our project, and attempted a demo with our client \\ \hline
		Week 2 & Random debug data was implemented, graph is now looking complete
           & Challenges in debug data and variation and frequency
           & Implementing a timer object and multiple random variables \\ \hline
		Week 3 & Positives & Deltas & Actions \\ \hline
		Week 4 & Positives & Deltas & Actions \\ \hline

    Winter Term & -- & -- & -- \\ \hline
    Week 1 & Positives & Deltas & Actions \\ \hline
		Week 2 & Positives & Deltas & Actions \\ \hline
		Week 3 & Other than UI, the application is working as planned
           & UI issues still persist, rotated axis titles are not sized properly
           & Research shows that rotated text doesn't allow auto-formatting \\ \hline
		Week 4 & Fonts were sized appropriately, main UI uses Views now
           & Even in Views, as opposed to Stacks, the rotating causes problems
           & Using restraints we were able to manage the rotated text slightly \\ \hline
    Week 5 & Positives & Deltas & Actions \\ \hline
    Week 6 & Positives & Deltas & Actions \\ \hline

	\end{tabular}
\end{center}

\newpage
\bibliographystyle{IEEEtran}
\bibliography{bibliography}
\end{document}
